\subsection{LCT (Link-Cut Tree)}
反正我不会，看着乐呵乐呵好了。。

这里求得是链 xor
\begin{lstlisting}
#include<bits/stdc++.h>
using namespace std;
#define ll int
#define f(i,a,b) for(ll i=a;i<=b;i++)
#define wt int tt=d;while(tt--)
#define py puts("Yes")
#define pn puts("No")
#define pritnf printf
#define fe(i,e) for(int i=0;i<e.size();i++)
#define vi vector<ll>
inline ll rd() {
	ll x=0,f=1;
	char c=getchar();
	while(!isdigit(c)){if(c=='-')f=-1;c=getchar();}
	while(isdigit(c))x=x*10+c-'0',c=getchar();
	return x*f;
}
ll dx[4]={0,1,0,-1};
ll dy[4]={1,0,-1,0};
#define d rd()
#define pb push_back
const ll N=300010;
struct edge{ll v,w,nx;}e[N<<1];
ll hd[N],cnt;
void add(ll u,ll v,ll w){e[++cnt]=(edge){v,w,hd[u]};hd[u]=cnt;}
ll qp(ll a,ll b,ll p){
	ll ans=1;while(b){
		if(b&1)ans=ans*a%p;
		a=a*a%p;b>>=1;
	}return ans;
}ll n,m;
ll Rt,tot;
ll fa[100010],ch[100010][2];
ll val[100010],sum[100010];
ll tg[100010];
#define ls(p) ch[p][0]
#define rs(p) ch[p][1]
bool isr(ll p){return (p!=ls(fa[p])&&p!=rs(fa[p]));}
void upd(ll p){sum[p]=sum[ls(p)]^sum[rs(p)]^val[p];}
bool get(ll p){return p==rs(fa[p]);}
void upt(ll p){swap(ls(p),rs(p));tg[p]^=1;}
void pd(ll p){
	if(!tg[p])return;
	if(ls(p))upt(ls(p));
	if(rs(p))upt(rs(p));
	tg[p]=0;
}void rot(ll x){
	ll y=fa[x],z=fa[y],t=get(x);
	if(!isr(y))ch[z][y==rs(z)]=x;
	ch[y][t]=ch[x][t^1];if(ch[x][t^1])fa[ch[x][t^1]]=y;
	ch[x][t^1]=y;fa[y]=x,fa[x]=z;upd(y);upd(x);
}void pda(ll x){if(!isr(x))pda(fa[x]);pd(x);}
void splay(ll x){//cout<<x<<endl;
	pda(x);//cout<<x<<endl;
	while(!isr(x)){
		ll y=fa[x],z=fa[y];
		if(!isr(y))rot(get(x)==get(y)?y:x);
		rot(x);
	}upd(x);
}void access(ll x){
	for(int y=0;x;y=x,x=fa[x]){
		splay(x);
		// cout<<x<<" "<<y<<endl;
		rs(x)=y,upd(x);
	}
}void mkr(ll x){
	access(x);
	splay(x);upt(x);
}ll find(ll x){
	access(x),splay(x);
	for(pd(x);ls(x);pd(x))x=ls(x);
	splay(x);return x;
}void split(ll x,ll y){
	mkr(x);access(y);splay(y);
}ll link(ll x,ll y){
	// cout<<x<<" "<<y<<endl;
	mkr(x);
	if(find(y)==x)return 0;
	fa[x]=y;splay(y);return 1;
}ll cut(ll x,ll y){
	split(x,y);
    if(ls(y)==x&&!rs(x)){ls(y)=fa[x]=0;upd(y);return 1;}return 0;
}
int main(){
	n=d,m=d;
    f(i,0,n)fa[i]=ch[i][0]=ch[i][1]=tg[i]=val[i]=sum[i]=0;
    f(i,1,n)sum[i]=val[i]=d;
	while(m--){
		ll o=d,x=d,y=d;
		if(o==0)split(x,y),printf("%d\n",sum[y]);
		if(o==1)link(x,y);if(o==2)cut(x,y);
		if(o==3)splay(x),val[x]=y,upd(x);
	}
	return 0;
}
\end{lstlisting}

\paragraph{多测与注意} main 的单组部分可放进 solve，初始化已补；仍是原来的数组版路径异或。查询默认两点连通，深链留意递归下传的栈空间。
